\chapter{Pengurangan: Langkah Pertama}\label{ch:g1:subtraction}

Tujuh balon, dua terbang pergi: berapa yang tersisa? Mengambil
sebagian berarti \emph{mengurangkan}, ditulis dengan tanda $-$. Dan
pengurangan punya wajah kedua: ia juga menjawab ``berapa lebihnya?''.

\section{Mengambil}

\begin{definition}[Pengurangan]\label{def:g1:subtraction:def}
\emph{Mengurangkan} berarti mengambil sebagian lalu membilang yang
tersisa. Kita menulis
\[
7 - 2 = 5
\]
dan dibaca: ``tujuh kurang dua sama dengan lima''.
\end{definition}

\begin{omfigure}
\begin{tikzpicture}[scale=0.6]
  \foreach \i in {0,...,4} \fill[omDef] (\i*1.05,0) circle (0.3);
  \foreach \i in {5,6} {
    \fill[omDef!30] (\i*1.05,0) circle (0.3);
    \draw[omThm, thick] (\i*1.05-0.32,-0.32) -- (\i*1.05+0.32,0.32);
    \draw[omThm, thick] (\i*1.05-0.32,0.32) -- (\i*1.05+0.32,-0.32);
  }
  \node[right, font=\small] at (7.6,0)
    {$7 - 2 = 5$: coretlah dua, tinggal lima};
\end{tikzpicture}

{\small Pengurangan dalam gambar: mencoret berarti mengambil.}
\end{omfigure}

\begin{example}[Membilang mundur]\label{ex:g1:subtraction:countback}
Untuk pengambilan kecil, bilanglah mundur pada garis bilangan:
$11 - 3$: sebutlah ``sebelas'', lalu tiga langkah mundur ---
``sepuluh, sembilan, delapan''. Jadi $11 - 3 = 8$.
\end{example}

\section{Berapa lebihnya?}

\begin{example}[Selisih]\label{ex:g1:subtraction:difference}
``Nina punya $9$ kelereng dan Tono punya $6$. Berapa lebihnya
kelereng Nina?'' Tidak ada yang diambil, tetapi ini tetap sebuah
pengurangan: pasangkanlah kelerengnya, lalu bilang kelereng Nina yang
tidak punya pasangan:
\[
9 - 6 = 3 .
\]
Nina punya $3$ lebih banyak. Hasil sebuah pengurangan disebut
\emph{selisih}\index{selisih}.
\end{example}

\begin{method}[Membilang maju]\label{met:g1:subtraction:countup}
Untuk mencari \omterm{ex:g1:subtraction:difference}{selisih}, membilang \emph{maju} sering kali paling
mudah: untuk $12 - 9$, mulailah dari $9$ lalu naik sampai $12$:
``sepuluh, sebelas, dua belas'' --- tiga langkah. Jadi $12 - 9 = 3$.
(Naiklah lewat $10$ untuk jarak yang lebih besar: $14 - 8$: dari $8$
ke $10$ ada $2$, dari $10$ ke $14$ ada $4$: \omterm{ex:g1:subtraction:difference}{selisihnya} $6$.)
\end{method}

\begin{omfigure}
\begin{tikzpicture}[scale=0.6]
  \draw[-Stealth] (5.6,0) -- (15.4,0);
  \foreach \i in {6,...,15} \draw (\i,-0.1) -- (\i,0.1);
  \foreach \i in {8,10,14}
    \node[below, font=\footnotesize] at (\i,-0.12) {\i};
  \draw[-Stealth, omDef, thick] (8,0.3) .. controls (9,0.9) ..
    (10,0.3);
  \node[omDef, font=\small] at (9,1.0) {$+2$};
  \draw[-Stealth, omThm, thick] (10,0.3) .. controls (12,1.2) ..
    (14,0.3);
  \node[omThm, font=\small] at (12,1.25) {$+4$};
\end{tikzpicture}

{\small $14 - 8$ dengan membilang maju: dari $8$ ke $14$ ada
$2 + 4 = 6$ langkah.}
\end{omfigure}

\section{Penjumlahan dan pengurangan adalah saudara kembar}

\begin{example}[Keluarga bilangan]\label{ex:g1:subtraction:family}
Bilangan $3$, $4$ dan $7$ membentuk satu keluarga kecil:
\[
3 + 4 = 7, \qquad
4 + 3 = 7, \qquad
7 - 4 = 3, \qquad
7 - 3 = 4 .
\]
Mengetahui satu fakta keluarga ini memberikan tiga sisanya secara
cuma-cuma. Pengurangan membatalkan penjumlahan: setelah $+4$,
melakukan $-4$ membawamu kembali.
\end{example}

\begin{example}[Memeriksa]\label{ex:g1:subtraction:check}
Sudah menghitung $13 - 5 = 8$? Periksalah dengan penjumlahan
kembarannya: $8 + 5 = 13$. \checkmark\ Sebuah pengurangan benar
tepat ketika penjumlahan baliknya cocok.
\end{example}

\begin{example}[Operasi yang mana?]\label{ex:g1:subtraction:choose}
``$6$ burung di kawat, $4$ terbang pergi'': mengambil, $6 - 4 = 2$.
``$6$ burung, lalu $4$ lagi datang'': menggabungkan, $6 + 4 = 10$.
``$6$ burung pipit dan $4$ burung kutilang: berapa lebih banyak
burung pipit?'': \omterm{ex:g1:subtraction:difference}{selisih}, $6 - 4 = 2$. Gambarlah ceritanya kalau ragu!
\end{example}

\section{Latihan}

\begin{exercise}[$\star$]\label{exo:g1:subtraction:1}
Gambar, coret, lalu hitunglah: $6 - 2$; \quad $8 - 5$; \quad
$10 - 4$.
\end{exercise}

\begin{exercise}[$\star$]\label{exo:g1:subtraction:2}
Hitunglah dengan membilang mundur: $9 - 3$; \quad $12 - 2$; \quad
$15 - 4$.
\end{exercise}

\begin{exercise}[$\star$]\label{exo:g1:subtraction:3}
Carilah \omterm{ex:g1:subtraction:difference}{selisihnya} dengan membilang maju: $11 - 8$; \quad $13 - 9$;
\quad $16 - 12$.
\end{exercise}

\begin{exercise}[$\star$]\label{exo:g1:subtraction:4}
Tulislah empat fakta dari keluarga $2$, $6$ dan $8$.
\end{exercise}

\begin{exercise}[$\star$]\label{exo:g1:subtraction:5}
Hitunglah, lalu periksa dengan penjumlahan kembarannya: $14 - 6$;
\quad $17 - 9$.
\end{exercise}

\begin{exercise}[$\star$]\label{exo:g1:subtraction:6}
Lengkapilah bagian yang kosong:
\[
10 - \;?\; = 7, \qquad
\;?\; - 5 = 5, \qquad
12 - \;?\; = 12 .
\]
\end{exercise}

\begin{exercise}[$\star$]\label{exo:g1:subtraction:7}
``Ada $15$ ceri di dalam mangkuk. Anak-anak memakan $8$. Berapa yang
tersisa?'' Tulislah pengurangannya dan kalimat jawabannya.
\end{exercise}

\begin{exercise}[$\star$]\label{exo:g1:subtraction:8}
Untuk setiap cerita, pilihlah $+$ atau $-$, lalu hitunglah: ``$9$
itik di kolam, $5$ berenang pergi''; ``$9$ itik, $5$ angsa datang'';
``Ali berumur $9$, adiknya berumur $5$: berapa tahun Ali lebih tua?''.
\end{exercise}

\begin{exercise}[$\star\star$]\label{exo:g1:subtraction:9}
Induk itik punya $12$ telur. Sebagian menetas; sekarang tersisa
$12 - ?$ telur dan ia membilang $7$ telur yang belum menetas. Berapa
telur yang menetas? (Ini keluarga bilangan yang mana?)

\end{exercise}
